% PAGES: 328-330

where $y = Bx$. Since $A$ is symmetric positive definite, it follows from (5.17)
that
\begin{align}
y^tAy &\ge 0, \quad \forall\, y \in \R^n; \\
y^tAy &= 0 \iff y = 0.
\end{align}
Then, from (10.126) and using (10.127--10.128), we obtain that
\begin{align}
x^tB^tABx &\ge 0, \quad \forall\, x \in \R^n; \\
x^tB^tABx &= 0 \iff y^tAy = 0 \iff y = 0 \iff Bx = 0 \iff x = 0,
\end{align}
where the last equivalence follows from the fact that $B$ is a nonsingular matrix.

From (10.129) and (10.130), we conclude that the matrix $B^tAB$ is symmetric
positive definite.
\end{proof}

\begin{lemma}
Let $A$ be an $n \times n$ symmetric positive matrix, and let $A_{n-1} =
A(2:n,2:n)$. Then, the matrix
\[
A_{n-1} \;-\; \frac{(A(1,2:n))^tA(1,2:n)}{A(1,1)}
\]
is symmetric positive definite.
\end{lemma}

\begin{proof}
Note that, since $A$ is a symmetric matrix, $A(2:n,1) = (A(1,2:n))^t$. Then, from
Lemma 10.15, we find that
\[
\resizebox{\linewidth}{!}{$\displaystyle
A \;=\; \begin{pmatrix} \sqrt{A(1,1)} & \bfzero^t \\[0.6em]
\dfrac{(A(1,2:n))^t}{\sqrt{A(1,1)}} & I_{n-1} \end{pmatrix}
\begin{pmatrix} 1 & \bfzero^t \\[0.6em]
\bfzero & A_{n-1} - \dfrac{(A(1,2:n))^tA(1,2:n)}{A(1,1)} \end{pmatrix}
\begin{pmatrix} \sqrt{A(1,1)} & \dfrac{A(1,2:n)}{\sqrt{A(1,1)}} \\[0.6em]
0 & I_{n-1} \end{pmatrix}$}
\]
which can be written as
\begin{equation}
A \;=\; M^t \begin{pmatrix} 1 & \bfzero^t \\[0.6em]
\bfzero & A_{n-1} - \dfrac{(A(1,2:n))^tA(1,2:n)}{A(1,1)} \end{pmatrix} M,
\end{equation}
where $M = \begin{pmatrix} \sqrt{A(1,1)} & \dfrac{A(1,2:n)}{\sqrt{A(1,1)}} \\[0.6em]
0 & I_{n-1} \end{pmatrix}$ and $\mathbf{0}$ denotes the column vector of size
$n-1$ with all entries equal to $0$. Note that $M$ is a nonsingular matrix, since
$M$ is upper triangular and $\det(M) = A(1,1) \ne 0$; cf. (10.4).

Multiply (10.131) to the left by $(M^t)^{-1}$ and to the right by $M^{-1}$ and
obtain that
\begin{equation}
(M^t)^{-1}AM^{-1} \;=\; \begin{pmatrix} 1 & \bfzero^t \\[0.6em]
\bfzero & A_{n-1} - \dfrac{(A(1,2:n))^tA(1,2:n)}{A(1,1)} \end{pmatrix}.
\end{equation}

Recall from Lemma 1.8 that $(M^t)^{-1} = (M^{-1})^t$. Let $B = M^{-1}$. Then,
(10.132) can be written as
\begin{equation}
B^tAB \;=\; \begin{pmatrix} 1 & \bfzero^t \\[0.6em]
\bfzero & A_{n-1} - \dfrac{(A(1,2:n))^tA(1,2:n)}{A(1,1)} \end{pmatrix}.
\end{equation}

Let $x = \begin{pmatrix} 0 \\ x_{n-1} \end{pmatrix} \in \R^n$, where $x_{n-1} \in
\R^{n-1}$. Using (10.133), we obtain that
\begin{align}
x^tB^tABx & \\
&= (0\ \ x_{n-1}^t) \begin{pmatrix} 1 & \bfzero^t \\[0.6em]
\bfzero & A_{n-1} - \dfrac{(A(1,2:n))^tA(1,2:n)}{A(1,1)} \end{pmatrix}
\begin{pmatrix} 0 \\ x_{n-1} \end{pmatrix} \notag \\
&= x_{n-1}^t\left(A_{n-1} - \frac{(A(1,2:n))^tA(1,2:n)}{A(1,1)}\right)x_{n-1}.
\end{align}

Since $A$ is a symmetric positive definite matrix, it follows from Lemma 10.16 that
the matrix $B^tAB$ is also symmetric positive definite. Thus,
\begin{equation}
x^tB^tABx \;\ge\; 0, \quad \forall\, x_n \in \R^n;
\end{equation}
\begin{equation}
x^tB^tABx \;=\; 0 \iff x = 0.
\end{equation}

Then, from (10.134--10.137), it follows that
\[
x_{n-1}^t\left(A_{n-1} - \frac{(A(1,2:n))^tA(1,2:n)}{A(1,1)}\right)x_{n-1} \;\ge\;
0, \quad \forall\, x_{n-1} \in \R^{n-1};
\]
\[
x_{n-1}^t\left(A_{n-1} - \frac{(A(1,2:n))^tA(1,2:n)}{A(1,1)}\right)x_{n-1} \;=\; 0
\iff x_{n-1} = 0,
\]
and we conclude that the matrix $A_{n-1} - \frac{(A(1,2:n))^tA(1,2:n)}{A(1,1)}$ is
symmetric positive definite.
\end{proof}

We include below a proof by induction of the fact that every symmetric positive
definite matrix has a Cholesky decomposition; see also Theorem 6.1.

\begin{theorem}
Any symmetric positive definite matrix has a Cholesky decomposition.
\end{theorem}

\begin{proof}
We give a proof by induction. Assume that any symmetric positive definite matrix of
size $n-1$ is diagonalizable. We will show that any symmetric positive definite
matrix of size $n$ is also diagonalizable.

Let $A$ be an $n \times n$ symmetric positive definite matrix. Then, $A(2:n,1) =
(A(1,2:n))^t$, and, from Lemma 10.15, we find that
\begin{equation}
\resizebox{0.92\linewidth}{!}{$\displaystyle
A \;=\; \begin{pmatrix} \sqrt{A(1,1)} & \bfzero^t \\[0.6em]
\dfrac{(A(1,2:n))^t}{\sqrt{A(1,1)}} & I_{n-1} \end{pmatrix}
\begin{pmatrix} 1 & \bfzero^t \\[0.6em]
\bfzero & A_{n-1} - \dfrac{(A(1,2:n))^tA(1,2:n)}{A(1,1)} \end{pmatrix}
\begin{pmatrix} \sqrt{A(1,1)} & \dfrac{A(1,2:n)}{\sqrt{A(1,1)}} \\[0.6em]
0 & I_{n-1} \end{pmatrix}$}
\end{equation}
where $I_{n-1}$ is the $(n-1)\times(n-1)$ identity matrix, and $\mathbf{0}$ is the
column vector of size $n-1$ with all entries equal to $0$.

Recall from Lemma 10.17 that $A_{n-1} - \frac{(A(1,2:n))^tA(1,2:n)}{A(1,1)}$ is an
$(n-1)\times(n-1)$ symmetric positive definite matrix. Then, from the induction
hypothesis, it follows that there exists an $(n-1)\times(n-1)$ upper triangular
matrix $U_{n-1}$ with positive entries on the main diagonal such that
\[
A_{n-1} \;-\; \frac{(A(1,2:n))^tA(1,2:n)}{A(1,1)} \;=\; U_{n-1}^tU_{n-1}.
\]

Thus,
\begin{align}
\begin{pmatrix} 1 & \bfzero^t \\[0.6em]
\bfzero & A_{n-1} - \dfrac{(A(1,2:n))^tA(1,2:n)}{A(1,1)} \end{pmatrix} & \\
&= \begin{pmatrix} 1 & \bfzero^t \\ \bfzero & U_{n-1}^tU_{n-1} \end{pmatrix} \notag
\\
&= \begin{pmatrix} 1 & \bfzero^t \\ \bfzero & U_{n-1}^t \end{pmatrix}
\begin{pmatrix} 1 & \bfzero^t \\ \bfzero & U_{n-1} \end{pmatrix}.
\end{align}

Note that
\begin{equation}
\begin{pmatrix} \sqrt{A(1,1)} & \bfzero^t \\[0.6em]
\dfrac{(A(1,2:n))^t}{\sqrt{A(1,1)}} & I_{n-1} \end{pmatrix}
\begin{pmatrix} 1 & \bfzero^t \\ \bfzero & U_{n-1}^t \end{pmatrix}
\;=\; \begin{pmatrix} \sqrt{A(1,1)} & \bfzero^t \\[0.6em]
\dfrac{(A(1,2:n))^t}{\sqrt{A(1,1)}} & U_{n-1}^t \end{pmatrix};
\end{equation}
\begin{equation}
\begin{pmatrix} 1 & \bfzero^t \\ \bfzero & U_{n-1} \end{pmatrix}
\begin{pmatrix} \sqrt{A(1,1)} & \dfrac{A(1,2:n)}{\sqrt{A(1,1)}} \\[0.6em]
0 & I_{n-1} \end{pmatrix}
\;=\; \begin{pmatrix} \sqrt{A(1,1)} & \dfrac{A(1,2:n)}{\sqrt{A(1,1)}} \\[0.6em]
0 & U_{n-1} \end{pmatrix}.
\end{equation}

From (10.138--10.142), we obtain that
\begin{align}
A &= \resizebox{0.85\linewidth}{!}{$\displaystyle
\begin{pmatrix} \sqrt{A(1,1)} & \bfzero^t \\[0.6em]
\dfrac{(A(1,2:n))^t}{\sqrt{A(1,1)}} & I_{n-1} \end{pmatrix}
\begin{pmatrix} 1 & \bfzero^t \\ \bfzero & U_{n-1}^t \end{pmatrix}
\begin{pmatrix} 1 & \bfzero^t \\ \bfzero & U_{n-1} \end{pmatrix}
\begin{pmatrix} \sqrt{A(1,1)} & \dfrac{A(1,2:n)}{\sqrt{A(1,1)}} \\[0.6em]
0 & I_{n-1} \end{pmatrix}$} \notag \\
&= \begin{pmatrix} \sqrt{A(1,1)} & \bfzero^t \\[0.6em]
\dfrac{(A(1,2:n))^t}{\sqrt{A(1,1)}} & U_{n-1}^t \end{pmatrix}
\begin{pmatrix} \sqrt{A(1,1)} & \dfrac{A(1,2:n)}{\sqrt{A(1,1)}} \\[0.6em]
0 & U_{n-1} \end{pmatrix} \notag \\
&= U_n^tU_n,
\end{align}
where $U_n = \begin{pmatrix} \sqrt{A(1,1)} & \dfrac{A(1,2:n)}{\sqrt{A(1,1)}} \\[0.6em]
0 & U_{n-1} \end{pmatrix}$ is an $n \times n$ upper triangular matrix. Note that the
main diagonal entries of $U_n$ are positive entries, since $U_{n-1}$ was assumed to
have positive entries on the main diagonal.

Thus, we showed that the matrix $A$ has the Cholesky decomposition $A = U_n^tU_n$.
We conclude by induction that any symmetric positive definite matrix has a Cholesky
decomposition.
\end{proof}

\begin{theorem}
If it exists, the LU decomposition without row pivoting of a matrix is unique.
\end{theorem}

\begin{proof}
We give a proof by contradiction. Assume that the $n \times n$ nonsingular matrix
$A$ has two LU decompositions without row pivoting, i.e., assume that
\[
A = L_1U_1 \quad \text{and} \quad A = L_2U_2,
\]
where $L_1$ and $L_2$ are lower triangular matrices with entries equal to 1 on the
main diagonal and $U_1$ and $U_2$ are nonsingular upper triangular matrices. Then,
\begin{equation}
L_1U_1 \;=\; L_2U_2.
\end{equation}
% Proof continues onto PDF page 331 (folio 315), the first page of the next file
% in this chunk -- the \begin{proof} above is intentionally left open here.
